\documentclass[10pt,a4paper]{amsart}
\usepackage[margin=19mm]{geometry}
\usepackage{amsmath,amssymb,mathtools}
\usepackage[scr=rsfs,cal=euler]{mathalfa}
\usepackage{xfrac}
\usepackage[unicode,hidelinks,bookmarks=false]{hyperref}
\newcommand{\ie}{i.e.\ }
\newcommand{\cf}{cf.\ }
\newcommand{\resp}{respectively\ }
\newcommand{\Subsection}{\subsection}
\providecommand{\nobreakdash}{\nobreak\hbox{-}}
\setlength{\emergencystretch}{2em}
\multlinegap0pt
\usepackage{amssymb}
\usepackage{mathtools}
\usepackage{tikz}
\usepackage{xcolor}
\usepackage[all]{xy}
\newtheorem{proposition}{Proposition}
\def  \Char     {\mathrm{char}}
\def  \GL        {\mathrm{GL}}
\def  \Stab     {\mathrm{Stab}}
\def  \cl       {\mathcal{L}}
\title{Hecke algebra of $\GL_n$ over a 2-dimensional local field}
\author{Xuecai Ma}
\date{Source: arXiv:2601.01610v1 (CC BY 4.0)}
\begin{document}
\maketitle

\noindent\emph{Source and licence.} Hecke algebra of $\GL_n$ over a 2-dimensional local field. Version arXiv:2601.01610v1 (CC BY 4.0), distributed by arXiv under the Creative Commons Attribution 4.0 International licence, http://creativecommons.org/licenses/by/4.0/, as recorded in the arXiv licence metadata for this exact version and shown on the abstract page https://arxiv.org/abs/2601.01610v1. Full licence notice: \texttt{licenses/arxiv-2601.01610-CC-BY-4.0.txt}.\par\smallskip

\noindent\textit{Editorial scope.} Counted unit: the article's proposition proposition (source0.tex lines 834--836). The article's complete original proof of it is delivered with it (source0.tex lines 839--882, one \texttt{proof} environment of 44 lines). No mathematics is written, repaired or re-derived: the statement and the proof are the article's own lines, in the article's own notation, and no retained line is substituted or re-flowed.  No further source-local context is needed: every cross-reference of every retained line resolves inside the two delivered spans.  Nothing was omitted from any retained span, so the package carries no omission record.  No retained line contains a citation, so the wrapper carries no bibliography and no bibliography entry.  No retained line contains a figure environment, an image-inclusion command or drawing code, so no image is delivered and the package declares no asset dependency.  Editorial formatting only, in addition: an AMS \texttt{amsart} preamble that restates the article's own declarations and macros used by the retained lines, namely the environments \texttt{proposition} on separate counters, together with the standard packages those lines need.  The retained environments print in this document as Proposition 1 in order of appearance, whereas the article prints the same items with the numbering of its own full document.  The complete native source of the article as distributed in the arXiv e-print for this exact version is bundled unchanged as \texttt{sources/192091/source0.tex}.
\begin{proposition}
 $\cl^M( \GL_n(F))$ is a measurable representation  of $\GL_n(F)$.
\end{proposition}

\begin{proof}
	By the translation invariant properties of the two-dimensional integral, we have $f(h^{-1}x) $ belongs to $\cl^M(\GL_n(F))$. To prove it is a measurable representation, we need to prove that 
$\mathrm{Stab}_f $ is  a measurable subgroup of $\GL_n(F)$.   Since $\cl^M(\GL_n(F))$ is generated by $g^{A, \Gamma}$, where $g$ are Bruhat-Schwartz functions on $\GL_n(E)$. We may suppose that $g = \Char^{\GL_n(E)}_K$ for a compact  subset $K$ of $\GL_n(E)$, and  $f = g^{A, \Gamma}$.  We then have
$$
g^{A, \Gamma} = \Char_V,   \quad   V=\left\{x \in \GL_n(F) |  x  \in  A(I_n+ t_2^{\Gamma} p^{-1}(K))  \right \},
$$
$$
g^{A, \Gamma} (h^{-1}x) = \Char_U,   \quad  U=\left\{x \in \GL_n(F) | h^{-1}x  \in  A(I_n+ t_2^{\Gamma} p^{-1}(K))  \right \}.
$$
By the assumption of $f$,
$$
\Stab_f = \left\{ h \in \GL_n(F)|  V=U \right\}.
$$
Since we have $\Gamma_{i,j} >0$, we get
$$
\Stab_f = \Stab_{A} \cap  \Stab_{p^{-1}K}.
$$
Since $K$ is a compact subset of $\GL_n(E)$,  for every $y \in K$, we have a compact open neighbourhood $K_yy \subset K$, where $K_y$ is a compact open subgroup of $\GL_n(E)$. Then 
$$
\{ K_yy| y \in K\}
$$ 
is an open cover of $K $, it has a finite  sub-cover $\{K_1y_1, \cdots,  K_sy_s \}$. Let
$$
H=  \cap_{1 \leq i \leq s} K_i.
$$
Then $K$ is left $H$ invariant. We claim that   $p^{-1}(K)$ is left $p^{-1}(H)$ invariant.

For  $w \in p^{-1}H$, and $z \in p^{-1}(K)$, we have
$$
p(wz)=p(w) p(z) \subset K,
$$
so $wz \subset p^{-1}(K)$.  On the other hand,  we have $\Stab_{A} =I_n$.  Thus  we have
\begin{enumerate}
	\item $\Stab_f=\{I_n\}$  if $A  \neq I_n$, and
	 \item  $\Stab_f =  I_n+ t_2^{\Gamma}p^{-1}(H)$ if $A= I_n$. 

\end{enumerate}
 We have $I_n$ is a measure $0$ subgroup, and 
\begin{eqnarray*}
	\int_{\GL_n(F)} \Char_{I_n+t_2^{\Gamma}p^{-1 }(H)}dx  &=& \int_{M_n (F)}  \Char_{I_n+t_2^{\Gamma}p^{-1}(H)} |\det x|^{-n}  dx  \\
	&=&\int_{M_n(F)} (\Char_H)^{0,0}(x) dx.  \\
 \end{eqnarray*}
We have $I_n+ t_2^{\Gamma}p^{-1}H$ is a measurable subset of $\GL_n(F)$,   clearly it is a subgroup of $\GL_n(F)$,  thus it is a measurable subgroup.
\end{proof}

\end{document}
